% TEMPLATE-MILC-v3.tex - plantilla maestra para todas las guias MILC v3
% Ver scripts/generadores/build-guia-milc.py para la lista completa de
% claves esperadas. El script valida que no queden placeholders sin
% reemplazar antes de escribir el archivo final.

\documentclass[11pt,letterpaper]{article}
\usepackage[margin=1.75cm]{geometry}
\usepackage[table]{xcolor}
\usepackage{fontspec}
\usepackage[spanish]{babel}
\usepackage{tikz}
\usepackage[most]{tcolorbox}
\usepackage{tabularx}
\usepackage{array}
\usepackage{fancyhdr}
\usepackage{titlesec}
\usepackage{ragged2e}
\usepackage{enumitem}
\usepackage{microtype}

\setmainfont{Helvetica Neue}
\setsansfont{Helvetica Neue}
\setlength{\parindent}{0pt}
\setlength{\parskip}{6pt}
\hyphenpenalty=200
\exhyphenpenalty=200
\pretolerance=400
\tolerance=2000
\setlength{\emergencystretch}{1em}
\renewcommand{\arraystretch}{1.25}
\setlist{leftmargin=5mm,itemsep=1mm,topsep=1mm}
\newcolumntype{Y}{>{\RaggedRight\arraybackslash}X}

\definecolor{milcMagenta}{HTML}{E6007E}
\definecolor{milcVino}{HTML}{5A0038}
\definecolor{milcTurquesa}{HTML}{0093A5}
\definecolor{milcVerde}{HTML}{58A923}
\definecolor{milcMostaza}{HTML}{E5B400}
\definecolor{milcOcre}{HTML}{B86600}
\definecolor{milcNegro}{HTML}{241816}
\definecolor{milcGris}{HTML}{F7F7F4}
\definecolor{milcLinea}{HTML}{E8E2DD}
\definecolor{milcGradePrimary}{HTML}{0066FF}
\definecolor{milcGradeDark}{HTML}{003D99}
\definecolor{milcGradeSoft}{HTML}{D6E8FF}
\definecolor{milcGradeAccent}{HTML}{CFE7FF}
\definecolor{milcGradeText}{HTML}{FFFFFF}
\color{milcNegro}

\pagestyle{fancy}
\fancyhf{}
\lhead{\small Tecnología e Informática · Grado 11}
\rhead{\small Guía 26 · MILC v3}
\cfoot{\small\thepage}
\renewcommand{\headrulewidth}{0pt}

\newtcolorbox{titlebox}[1]{
  enhanced,colback=#1,colframe=#1,arc=7mm,boxrule=0pt,
  left=7mm,right=7mm,top=3mm,bottom=3mm,
  before skip=8pt,after skip=8pt,
  fontupper=\sffamily\bfseries\Large\color{white}
}

\newtcolorbox{softbox}[3]{
  enhanced,colback=#2,colframe=#1,borderline west={4pt}{0pt}{#1},
  arc=2mm,boxrule=.4pt,left=4mm,right=4mm,top=3mm,bottom=3mm,
  before skip=6pt,after skip=8pt,title={#3},coltitle=white,
  colbacktitle=#1,fonttitle=\sffamily\bfseries,fontupper=\RaggedRight,
  attach boxed title to top left={xshift=3mm,yshift=-2mm},
  boxed title style={arc=2mm, boxrule=0pt}
}

\newtcolorbox{drawbox}[2]{
  enhanced,colback=white,colframe=#1,arc=4mm,boxrule=1pt,
  left=5mm,right=5mm,top=4mm,bottom=4mm,before skip=8pt,after skip=8pt,
  title={#2},coltitle=white,colbacktitle=#1,fonttitle=\sffamily\bfseries,
  fontupper=\RaggedRight,
  attach boxed title to top left={xshift=4mm,yshift=-2mm},
  boxed title style={arc=3mm, boxrule=0pt}
}

\newtcolorbox{infoband}[1]{
  enhanced,colback=#1!8,colframe=#1!50,arc=2mm,boxrule=.5pt,
  left=4mm,right=4mm,top=2mm,bottom=2mm,before skip=4pt,after skip=4pt,
  fontupper=\small\RaggedRight
}

% Puente narrativo entre secciones (contrato editorial Conéctate).
% Da continuidad: "Acabas de X. Ahora vas a Y porque Z."
% Visualmente: banda izquierda en color del grado, texto en itálica pequeña.
\newtcolorbox{puentebox}{
  enhanced,colback=milcGradeSoft,colframe=milcGradeDark,
  borderline west={3pt}{0pt}{milcGradeDark},
  arc=0mm,boxrule=0pt,left=5mm,right=4mm,top=2.5mm,bottom=2.5mm,
  before skip=4pt,after skip=8pt,
  fontupper=\itshape\small\color{milcNegro}\RaggedRight
}
\newcommand{\puente}[1]{%
  \begin{puentebox}%
    \textcolor{milcGradeDark}{\textbf{\textrightarrow}}\hspace{2mm}#1%
  \end{puentebox}%
}

\newcommand{\linead}[1]{\noindent\color{milcLinea}\rule{#1}{0.5pt}\color{milcNegro}}
\newcommand{\respuesta}[1]{\par\vspace{2mm}\foreach \n in {1,...,#1}{\noindent\color{milcLinea}\rule{\linewidth}{0.5pt}\par\vspace{5mm}}\color{milcNegro}}
\newcommand{\checkbox}{\textcolor{milcGradeDark}{\fbox{\phantom{x}}}\hspace{5pt}}

\newcommand{\phasecard}[4]{
\begin{tcolorbox}[enhanced,colback=#3,colframe=#2,arc=3mm,boxrule=.5pt,height=3.05cm,valign=center,halign=center,left=1.5mm,right=1.5mm,top=2mm,bottom=2mm]
{\sffamily\bfseries\fontsize{22}{24}\selectfont\textcolor{#2}{#1}}\par
{\sffamily\bfseries\small #4}
\end{tcolorbox}
}

\begin{document}
\thispagestyle{empty}
\begin{tikzpicture}[remember picture,overlay]
  \fill[white] (current page.south west) rectangle (current page.north east);
  \path[fill=milcGradePrimary,rounded corners=16mm]
    ([xshift=.55cm,yshift=-.55cm]current page.north west) rectangle
    ([xshift=-.55cm,yshift=3.65cm]current page.south east);

  \node[anchor=north west,text=milcGradeText] at ([xshift=2.0cm,yshift=-1.85cm]current page.north west)
    {\sffamily\bfseries\fontsize{14}{16}\selectfont GRADO};
  \node[anchor=north west,text=milcGradeText,opacity=.82] at ([xshift=2.0cm,yshift=-2.75cm]current page.north west)
    {\sffamily\bfseries\fontsize{9}{11}\selectfont SERIE GUÍAS · TECNOLOGÍA E INFORMÁTICA};

  \begin{scope}[shift={([xshift=-3.05cm,yshift=-2.55cm]current page.north east)}]
    \fill[white,opacity=.80] (-.62,-.52) rectangle (.62,.52);
    \draw[milcGradeText,line width=1pt] (-.62,-.52) rectangle (.62,.52);
    \fill[milcGradeDark] (-.42,-.38) rectangle (-.22,.06);
    \fill[milcGradeAccent] (-.10,-.38) rectangle (.10,.24);
    \fill[milcGradePrimary!60!black] (.22,-.38) rectangle (.42,.38);
  \end{scope}

  \node[anchor=west,text=milcGradeText] at ([xshift=1.9cm,yshift=-12.85cm]current page.north west)
    {\sffamily\bfseries\fontsize{124}{124}\selectfont 11};
  \draw[milcGradeText,line width=1.7pt]
    ([xshift=4.52cm,yshift=-11.68cm]current page.north west) circle (.13cm);

  \node[anchor=west,text=milcGradeText,text width=8.7cm,align=left] at ([xshift=10.35cm,yshift=-11.6cm]current page.north west)
    {
     {\sffamily\bfseries\fontsize{19}{22}\selectfont Pitch ---\\contar la historia\\de tu proyecto}\\[4mm]
     {\sffamily\bfseries\fontsize{23}{26}\selectfont Guía 26-11-TIC}\\[2mm]
     {\sffamily\fontsize{13}{16}\selectfont Período 3 · Proyecto final emprendedor}\\[5mm]
     {\sffamily\bfseries\fontsize{11}{13}\selectfont Institución Educativa Sor María Juliana}\\[1mm]
     {\sffamily\fontsize{10}{12}\selectfont Autor: Dr. Álvaro Cárdenas Orozco}
    };

  \path[fill=milcGradeSoft,rounded corners=7mm]
    ([xshift=1.35cm,yshift=.85cm]current page.south west) rectangle
    ([xshift=-1.35cm,yshift=3.15cm]current page.south east);
  \draw[milcGradeDark,line width=.65pt,rounded corners=7mm]
    ([xshift=1.35cm,yshift=.85cm]current page.south west) rectangle
    ([xshift=-1.35cm,yshift=3.15cm]current page.south east);
  \node[anchor=center,text=milcNegro,text width=17.0cm,align=center] at ([yshift=2.0cm]current page.south)
    {\sffamily\fontsize{10}{12}\selectfont Metodología MILC · Apertura ancestral · Pensamiento computacional · Triángulo Dussel-Estoicismo-Floridi\\
    \textcolor{milcGradeDark}{\bfseries Producto:} Pitch deck de 7 slides estructurado (problema, afectados, solución, MVP, datos, modelo, llamada) + guion escrito de 3 minutos con tiempos por slide + grabación de ensayo de 3 minutos exactos hecha con celular, escuchada y autocorregida al menos 2 veces};
\end{tikzpicture}
\mbox{}
\newpage

\section*{Apertura: Pitch --- contar la historia de tu proyecto}

\begin{softbox}{milcGradeDark}{milcGradeSoft}{Saber ancestral}
En los corregimientos del Valle del Cauca, en las veredas del Chocó, en las noches de fogón del Pacífico, hay un oficio que sostuvo la educación cuando no había escuelas: \textbf{contar el cuento}. La abuela que enseñaba a los nietos por qué no se camina solo de noche; el abuelo que explicaba el origen del nombre del pueblo; la mayora que transmitía la receta de la viuda de pescado con un relato y no con una lista de ingredientes. El cuento ancestral tenía una arquitectura silenciosa que toda comunidad oral conocía sin nombrarla: (1) \textbf{un comienzo que enganchaba} con una imagen concreta o una pregunta; (2) \textbf{un conflicto que dolía} (alguien quería algo, algo se interpuso); (3) \textbf{una resolución que dejaba enseñanza}. Aristóteles llamó a esto \emph{principio, medio y fin}; los abuelos del Valle lo conocían siglos antes de que el filósofo griego lo escribiera. Quien sabía contar bien era escuchado, recordado, replicado; quien no sabía, hablaba al vacío. \textbf{La narrativa es estructura ancestral, no truco moderno de marketing.}
\end{softbox}

\begin{softbox}{milcTurquesa}{milcGris}{Saber contemporáneo}
El \textbf{pitch} (en inglés, \emph{lanzamiento}) es la forma contemporánea más exigente de la narrativa: contar un proyecto entero en 3-5 minutos con estructura narrativa, evidencia y llamada a la acción. Los pitch profesionales más reconocidos (Y Combinator, Sequoia, Bogotá Pitch Fest) comparten una arquitectura de \textbf{7 slides estándar}: (1) Problema, (2) Afectados, (3) Solución, (4) MVP/Producto, (5) Datos/Tracción, (6) Modelo/Negocio, (7) Llamada/Pedido. La regla profesional famosa: \textbf{si no puedes contar tu proyecto en 3 minutos, todavía no lo entiendes}. El pitch no es presentación: es \textbf{filtro narrativo}. Te obliga a destilar 10 sesiones de trabajo en su esencia narrativa más pura, con un comienzo que enganche, un conflicto que duela, una resolución que inspire y una acción concreta que pidas a quien escucha (inversión, alianza, retroalimentación, voto, contratación). El pitch es donde se compromete o se pierde el respeto profesional.
\end{softbox}

\begin{softbox}{milcMostaza}{milcGris}{Pregunta-puente}
\textit{¿Qué sabía la abuela del Valle al contar el cuento del fogón, que el emprendedor novato olvida cuando llena 20 slides con tablas y gráficos? ¿Y por qué los pitch más exitosos del mundo se parecen estructuralmente más a un cuento de fogón que a una clase universitaria?}
\end{softbox}

\begin{softbox}{milcVerde}{milcGris}{Saber y saber hacer}
Al terminar podrás: (1) \textbf{identificar} la \emph{historia mínima} de tu proyecto, capaz de sostener un pitch de 3 minutos: el problema que duele, la solución que alivia, la evidencia que prueba; (2) \textbf{analizar} la arquitectura de 7 slides estándar y aplicarla a tu proyecto con honestidad (no inflando ni inventando); (3) \textbf{crear} tu pitch deck de 7 slides + guion escrito + grabación de ensayo de 3 minutos exactos; (4) \textbf{evaluar} tu pitch escuchándolo grabado, identificando 3 debilidades concretas (verbosidad, rigidez, miedo, falta de evidencia) y corrigiendo al menos 2.
\end{softbox}

\small
\begin{tabularx}{\textwidth}{>{\bfseries}p{3.0cm}Y}
\rowcolor{milcGradePrimary!12}Autor & Dr. Álvaro Cárdenas Orozco\\
DBA & Diseño, implementación y sustentación de un proyecto de impacto local (MEN, T\&I)\\
Referentes & Dussel (analéctica · sur global) · Lean Startup · Floridi · Estoicismo\\
Producto final & Pitch deck de 7 slides estructurado (problema, afectados, solución, MVP, datos, modelo, llamada) + guion escrito de 3 minutos con tiempos por slide + grabación de ensayo de 3 minutos exactos hecha con celular, escuchada y autocorregida al menos 2 veces\\
\end{tabularx}
\normalsize

\puente{Antes de empezar, mira el viaje completo. Llevas 5 sesiones de trabajo silencioso: problema (S1), validación (S2), modelo (S3), MVP (S4), datos (S5). Hoy todo eso \textbf{se convierte en historia} de 3 minutos. La sesión 7 trabajará marca, la 8 finanzas; pero el pitch de hoy es la primera vez que el proyecto se vuelve público. Las próximas sesiones afinan; esta sesión define el tono narrativo del cierre.}

\begin{titlebox}{milcGradeDark}
Ruta MILC: así vamos a aprender
\end{titlebox}
\begin{tabularx}{\textwidth}{>{\centering\arraybackslash}X>{\centering\arraybackslash}X>{\centering\arraybackslash}X>{\centering\arraybackslash}X}
\phasecard{1}{milcVerde}{milcGris}{Escuta\\[-1mm]\small Observo, escucho y pregunto.} &
\phasecard{2}{milcTurquesa}{milcGris}{Sistematización\\[-1mm]\small Pensamiento computacional.} &
\phasecard{3}{milcMagenta}{milcGris}{Praxis\\[-1mm]\small Construyo y aplico.} &
\phasecard{4}{milcVino}{milcGris}{Evaluación\\[-1mm]\small Reflexión liberadora.}
\end{tabularx}


\puente{Antes de teorizar, vas a hacer una \emph{destilación}: en 3 minutos personales, cuéntale a un compañero o a ti mismo en voz alta lo que has hecho en las sesiones 1-5. Sin guion, sin slides, sin filtros. Esa primera versión cruda revela qué partes naturalmente recuerdas como importantes y qué partes están floja o ausente.}

\begin{titlebox}{milcVerde}
Fase 1 · Escuta
\end{titlebox}

\begin{softbox}{milcVerde}{milcGris}{Escena para escuchar}
\textbf{Actividad 1 · IDENTIFICA --- Cuento crudo de mi proyecto en 3 minutos} (15 min · individual). Sin slides, sin guion escrito, sin filtros: cuéntate a ti mismo (en voz alta, frente al cuaderno o a un espejo, o grabándote con el celular sin editar) lo que has hecho en las sesiones 1 a 5. Reglas: (1) máximo 3 minutos cronometrados; (2) cuenta como si le hablaras a un familiar o amigo, no a un profesor; (3) no leas de la guía: \emph{recuerda}. Al terminar, anota: ¿qué naturalmente recordé primero?, ¿qué olvidé?, ¿qué parte sonó floja o vacía? Esta primera versión cruda es la materia prima del pitch.
\end{softbox}

\checkbox Conté el proyecto en voz alta sin leer guion.\\
\checkbox Me ajusté a 3 minutos cronometrados.\\
\checkbox Anoté qué recordé primero, qué olvidé, qué sonó flojo.

\begin{infoband}{milcVerde}
\textbf{Lo que va al cuaderno (Actividad 1 · IDENTIFICA).} Título: «Actividad 1 --- Cuento crudo de mi proyecto». \textbf{Formato:} 3 párrafos: (a) qué recordé naturalmente primero, (b) qué olvidé que probablemente importa, (c) qué parte sonó floja o vacía. \textbf{Sabes que terminaste cuando} sabes cuáles son tus 2-3 partes débiles narrativamente.
\end{infoband}

\puente{Ya tienes la versión cruda. Ahora vas a entender la \textbf{anatomía del pitch profesional}: los 7 slides estándar con su contenido específico, la regla de tiempo (1 slide ≈ 25-30 segundos), los 6 errores típicos del pitch novato (slide repleta de texto, hablar de uno y no del usuario, datos sin baseline, sin pedido al final, voz monótona, exceso de detalle técnico).}

\begin{titlebox}{milcTurquesa}
Fase 2 · Sistematización con pensamiento computacional
\end{titlebox}

\begin{softbox}{milcTurquesa}{milcGris}{Cuatro pilares aplicados al tema}
Un pitch novato es 20 slides desordenadas. Un pitch profesional es 7 slides con arquitectura narrativa. Para que el tuyo sea del segundo tipo necesitas la anatomía estándar del Y Combinator y los errores típicos que matan la atención del oyente en el segundo 20.
\end{softbox}

\small
\begin{tabularx}{\textwidth}{>{\bfseries\RaggedRight\arraybackslash}p{3.6cm}Y}
\rowcolor{milcTurquesa!90}\color{white}Pilar & \color{white}\bfseries Aplicación al tema\\
1. Descomposición & El pitch profesional se descompone en 7 slides: (1) \textbf{Problema} (25 seg): el dolor real, con persona concreta y dato (sesión 1-2). (2) \textbf{Afectados} (20 seg): cuántas personas viven este dolor y dónde están. (3) \textbf{Solución} (25 seg): qué propones, en 1 frase con verbo de cambio. (4) \textbf{MVP/Producto} (30 seg): imagen del MVP de la sesión 4 + las 3 piezas (captura, valor, respuesta). (5) \textbf{Datos/Tracción} (30 seg): los 3 hallazgos numéricos de la sesión 5 con honestidad de muestra. (6) \textbf{Modelo} (25 seg): cómo gana plata el proyecto (Canvas reducido a 1 línea). (7) \textbf{Llamada/Pedido} (25 seg): qué le pides al oyente concretamente (alianza, retroalimentación, contacto, inversión, voto, contratación).\\
2. Reconocimiento de patrones & La \textbf{regla de tiempo} es ley de oficio: 7 slides × 25-30 segundos = 3 a 3,5 minutos. Si una slide te toma 1 minuto, otras dos quedaron sin tiempo. La phronesis del pitch consiste en \textbf{distribuir el tiempo equitativamente}, con 5 segundos extra solo para la slide de datos (la más importante junto con el problema). El cronómetro es tu mejor crítico.\\
3. Abstracción & Lo esencial cabe en una pregunta: \emph{¿el oyente se interesará por el problema antes del segundo 20?}. Si la respuesta es no, perdiste el pitch entero. La slide 1 (problema) es la más importante: tiene que enganchar con imagen concreta, persona real, dato medible, no con generalidades. Los pitch que arrancan con \emph{``en Colombia hay X millones de…''} pierden a la audiencia antes de la primera idea.\\
4. Algoritmo conceptual & Algoritmo: (1) escribe el guion de 3 minutos en prosa antes de hacer las slides (el guion manda, las slides ilustran); (2) descompón el guion en 7 párrafos correspondientes a los 7 slides; (3) reduce cada slide a \emph{1 idea visual + máximo 7 palabras} (regla 7×7 invertida: no más de 7 palabras por slide); (4) diseña con consistencia visual mínima (1 tipografía, 2-3 colores, espacios respirables); (5) ensaya en voz alta sin público; (6) graba con celular y escúchate; (7) corrige 2 debilidades concretas (verbosidad, voz monótona, rigidez, miedo, falta de evidencia).\\
\end{tabularx}
\normalsize

\begin{drawbox}{milcTurquesa}{Anatomía del pitch de 7 slides}
\textbf{1. Problema (25 seg):} dolor con persona y dato. \\[1mm] \textbf{2. Afectados (20 seg):} cuántos y dónde. \\[1mm] \textbf{3. Solución (25 seg):} qué propones en 1 frase. \\[1mm] \textbf{4. MVP (30 seg):} imagen + 3 piezas. \\[1mm] \textbf{5. Datos (30 seg):} 3 hallazgos numéricos. \\[1mm] \textbf{6. Modelo (25 seg):} cómo gana plata. \\[1mm] \textbf{7. Pedido (25 seg):} qué le pides al oyente.
\end{drawbox}

\begin{softbox}{milcMostaza}{milcGris}{Errores comunes y su corrección}
(1) \textbf{Slide repleta de texto}: el oyente lee o escucha, no las dos; si lees lo de la slide, no estás pitcheando. (2) \textbf{Hablar de uno y no del usuario}: \emph{``mi proyecto, mi idea, yo creo''} en lugar de \emph{``María, la tendera, pierde…''}. (3) \textbf{Datos sin baseline ni contexto}: decir \emph{``5 inscritos''} sin comparar con meta ni muestra. (4) \textbf{Sin pedido al final}: cerrar con \emph{``muchas gracias''} en lugar de \emph{``los invito a…''}. (5) \textbf{Voz monótona o leyendo}: leer mata el pitch incluso con datos buenos. (6) \textbf{Exceso de detalle técnico}: explicar la arquitectura Zapier-Sheets-ManyChat a una audiencia que solo quería saber qué hace para el usuario. (7) \textbf{Pasarse del tiempo}: pasar de 3,5 min indica falta de preparación, no de contenido.
\end{softbox}

\begin{infoband}{milcTurquesa}
\textbf{Lo que va al cuaderno (Actividad 2 · ANALIZA).} Título: «Actividad 2 --- Anatomía del pitch». \textbf{Formato:} (a) plantilla de 7 slides con título y duración; (b) los 7 errores típicos con marca del que más temes cometer; (c) el guion de 3 minutos en prosa antes de hacer slides. \textbf{Sabes que terminaste cuando} podrías delegar el diseño visual del deck a otro compañero con solo ese guion.
\end{infoband}

\puente{Tienes el método. Toca aplicarlo: construir las 7 slides con frases cortas y 1 idea visual por slide, escribir el guion hablado de 3 minutos, ensayar en voz alta sin público, grabarte con el celular, escucharte 2 veces y corregir 2 debilidades.}

\begin{titlebox}{milcMagenta}
Fase 3 · Praxis con pensamiento computacional
\end{titlebox}

\begin{softbox}{milcMagenta}{milcGris}{Construyendo el producto con los cuatro pilares}
\textbf{Actividad 3 · CREA + EVALÚA --- Pitch deck + grabación + autocrítica} (45 min en sesión, 1-2 horas de ensayo en casa · individual). Hora de armar. Construyes el deck de 7 slides en Canva, Google Slides o Keynote; escribes el guion de 3 minutos en prosa; ensayas frente al espejo; grabas con celular; te escuchas 2 veces; identificas 3 debilidades; corriges 2.
\end{softbox}

\small
\begin{tabularx}{\textwidth}{>{\bfseries\RaggedRight\arraybackslash}p{3.6cm}Y}
\rowcolor{milcMagenta!90}\color{white}Pilar & \color{white}\bfseries Aplicación al producto\\
Descomposición & Tu pitch deck se construye en 6 pasos: (1) escribir el guion completo en prosa (700-800 palabras = 3 minutos hablados); (2) dividir en 7 párrafos correspondientes a slides; (3) diseñar cada slide con la regla de 7 palabras o menos + 1 imagen o gráfico; (4) usar plantilla mínima consistente (1 tipografía, 2-3 colores, espacios respirables); (5) ensayar 3 veces en voz alta sin grabar para que el cuerpo aprenda; (6) grabar con celular en horizontal, escucharse completo, corregir.\\
Patrones & Sigue el patrón \textbf{``palabra hablada sobre slide leída''}: las slides son apoyo visual, no transcripción. Si el oyente puede leer tu slide y entender, no necesita escucharte. La phronesis del pitch consiste en que la slide diga lo mínimo y tu voz diga lo importante. Steve Jobs llamó a esto \emph{``demand attention, not divide it''}.\\
Abstracción & La autocrítica al escuchar la grabación se estructura en 4 dimensiones: (a) \textbf{Tiempo}: ¿pasé de 3,5 min? ¿alguna slide se llevó demasiado tiempo? (b) \textbf{Voz}: ¿soné monótono o vivo? ¿leí o conté? (c) \textbf{Estructura}: ¿el oyente entendió el problema antes del segundo 20? ¿hice el pedido al final? (d) \textbf{Evidencia}: ¿los 3 hallazgos numéricos quedaron claros con contexto?\\
Algoritmo de armado & Algoritmo: (1) escribe el guion; (2) construye las slides; (3) ensaya 3 veces sin grabar; (4) graba; (5) escúchate completo sin pausar; (6) escúchate de nuevo anotando las 3 debilidades; (7) corrige las 2 más graves antes de la próxima sesión; (8) opcional: pide a un familiar que vea el pitch y te diga si entendió el problema.\\
\end{tabularx}
\normalsize

\begin{drawbox}{milcMagenta}{Checklist antes de entregar el pitch}
\checkbox 7 slides construidas con título y máximo 7 palabras de texto. \\[1mm] \checkbox Guion escrito en prosa de 700-800 palabras. \\[1mm] \checkbox Estructura: Problema, Afectados, Solución, MVP, Datos, Modelo, Pedido. \\[1mm] \checkbox Ensayado 3 veces en voz alta sin grabar. \\[1mm] \checkbox Grabación de 3 minutos exactos hecha con celular. \\[1mm] \checkbox Escucharon la grabación al menos 2 veces. \\[1mm] \checkbox 3 debilidades identificadas, 2 corregidas.
\end{drawbox}

\begin{softbox}{milcMagenta}{milcGris}{Plantilla de trabajo}
\textbf{Slide 1 --- Problema (25 seg).} \textit{Persona + dolor + dato.} \\[1mm] \textbf{Slide 2 --- Afectados (20 seg).} \textit{Cuántos y dónde.} \\[1mm] \textbf{Slide 3 --- Solución (25 seg).} \textit{Una frase con verbo de cambio.} \\[1mm] \textbf{Slide 4 --- MVP (30 seg).} \textit{Imagen + las 3 piezas.} \\[1mm] \textbf{Slide 5 --- Datos (30 seg).} \textit{Los 3 hallazgos numéricos.} \\[1mm] \textbf{Slide 6 --- Modelo (25 seg).} \textit{Cómo gana plata.} \\[1mm] \textbf{Slide 7 --- Pedido (25 seg).} \textit{Qué le pides al oyente.}
\end{softbox}

\begin{infoband}{milcMagenta}
\textbf{Lo que va al cuaderno (Actividad 3 · CREA + EVALÚA).} Título: «Actividad 3 --- Mi pitch». \textbf{Formato:} (a) guion completo escrito (700-800 palabras); (b) miniaturas de las 7 slides dibujadas o pegadas; (c) ficha de autocrítica con 3 debilidades identificadas y 2 corregidas; (d) link de la grabación o referencia al archivo. \textbf{Sabes que terminaste cuando} la grabación corregida dura entre 3 y 3,5 minutos exactos.
\end{infoband}

\puente{Lo que sigue resume \emph{qué} entregas hoy: pitch deck de 7 slides + guion escrito + grabación de ensayo de 3 minutos. El criterio principal: que un familiar que no conoce tu proyecto, después de ver tu pitch, pueda contarlo en sus propias palabras a otra persona en 1 minuto. Si no puede, el pitch falló su misión.}

\begin{titlebox}{milcMostaza}
Producto de la sesión
\end{titlebox}
\textbf{Pitch deck de 7 slides + guion escrito + grabación de ensayo de 3 min}.
\begin{softbox}{milcMostaza}{milcGris}{Criterios de calidad}
(1) Estructura de 7 slides estándar respetada. (2) Cada slide con máximo 7 palabras + 1 elemento visual. (3) Guion escrito en prosa de 700-800 palabras. (4) Grabación de 3 minutos exactos hecha con celular. (5) Autocrítica con 3 debilidades identificadas. (6) Al menos 2 debilidades corregidas antes de cerrar.
\end{softbox}

\puente{Antes de cerrar, mira el pitch desde las cinco dimensiones humanas. Contar tu proyecto con honestidad narrativa es respeto por la audiencia entera; inflar los datos para impresionar es traición al oficio.}

\begin{titlebox}{milcVino}
Fase 4 · Evaluación liberadora — cinco dimensiones
\end{titlebox}

\begin{softbox}{milcVino}{milcGris}{Sentido de la evaluación}
Evaluar no es solo revisar una nota. Es reconocer qué aprendí, qué cambió en mí, qué emociones manejé, cómo me relaciono con la ciudad y la comunidad, y qué le dejaré a quien viene detrás.
\end{softbox}

\textbf{Mi reflexión liberadora en cinco dimensiones.} Completa cada frase iniciada:

\small
\begin{tabularx}{\textwidth}{>{\bfseries\RaggedRight\arraybackslash}p{3.4cm}Y>{\RaggedRight\arraybackslash}p{4.6cm}}
\rowcolor{milcVino}\color{white}Dimensión & \color{white}\bfseries Mi respuesta (frame starter) & \color{white}\bfseries Algo que descubrí\\
1. Personal & Lo que más cambió en mí fue \linead{6.5cm} & \linead{4.2cm}\\
2. Emocional & La emoción más fuerte fue \linead{2.5cm} y la regulé \linead{3cm} & \linead{4.2cm}\\
3. Ciudadana & En democracia esto importa porque \linead{6cm} & \linead{4.2cm}\\
4. Local (Cartago) & En mi barrio sería útil para \linead{6.5cm} & \linead{4.2cm}\\
5. Intergeneracional & A mi abuela/abuelo le diría \linead{6.5cm} & \linead{4.2cm}\\
\end{tabularx}
\normalsize

\begin{infoband}{milcVino}
\textbf{Para la próxima sustentación.} Vuelve a esta tabla en seis meses. Verás que las respuestas cambian — no porque lo aprendido cambie, sino porque tú cambias. La evaluación liberadora es una conversación contigo a través del tiempo.
\end{infoband}


\puente{Y para terminar, tres voces sobre la palabra y la verdad. Dussel sobre el pitch que da voz al afectado o se la quita, Séneca sobre la sobriedad narrativa, Floridi sobre la ética de comunicar evidencia en la era de la atención escasa.}

\begin{titlebox}{milcGradeDark}
Triángulo de pensamiento · cierre filosófico
\end{titlebox}

\begin{softbox}{milcVino}{milcGris}{Dussel · lente del nosotros}
\textit{``Todo pitch decide si da voz al afectado o se la quita; no hay narrativa neutra.''}

Tu pitch puede tener al afectado en el centro (\emph{``María, la tendera del barrio Obrero, pierde 30 mil pesos a la semana porque…''}) o puede tenerse a sí mismo en el centro (\emph{``mi idea es genial porque…''}). La filosofía de la liberación pide lo primero: que la persona que sufre el problema aparezca con nombre, oficio y voz en la slide 1. Sin afectado nombrado, el pitch es propaganda. Con afectado nombrado, el pitch es honesto.

\textbf{En mi slide 1, ¿está el afectado con nombre y voz, o estoy yo con \emph{``mi proyecto''} como protagonista?}
\end{softbox}
\linead{\linewidth}\\[1mm]
\linead{\linewidth}

\begin{softbox}{milcMostaza}{milcGris}{Estoicismo · Séneca · cuidado interior}
\textit{``Habla sobrio o no hables; el discurso inflado revela carácter inflado.''}

Es tentador inflar los datos (\emph{``500 visitas''} cuando fueron 50), exagerar el impacto (\emph{``revolucionario''} cuando es solo útil), prometer escala (\emph{``escalable a millones''} cuando aún no tienes 10 usuarios). La sobriedad estoica recomienda lo contrario: \emph{di lo que es, ni más grande ni más pequeño}. El oyente experimentado detecta la inflación en el segundo 30 y deja de escuchar. La sobriedad es la mayor demostración de seguridad profesional.

\textbf{En mi pitch, ¿estoy diciendo los datos como son (n=10, decisión preliminar), o los estoy inflando para sonar más exitoso?}
\end{softbox}
\linead{\linewidth}\\[1mm]
\linead{\linewidth}

\begin{softbox}{milcTurquesa}{milcGris}{Floridi · lente de la infoesfera}
\textit{``Comunicar evidencia con honestidad es la nueva ética profesional en la era de la atención escasa.''}

La era de la información compite por atención: cada oyente tiene 100 pitches por mes pidiendo sus 3 minutos. La ética digital contemporánea exige no engañar para retener atención. Pitches que inflan, exageran o ocultan limitaciones erosionan la confianza colectiva en el ecosistema emprendedor. Tu pitch, aunque sea escolar, es entrenamiento ético: si te acostumbras a inflar hoy, lo harás cuando los montos sean reales.

\textbf{¿Mi pitch declara las limitaciones (muestra pequeña, etapa temprana) o las esconde para sonar más maduro?}
\end{softbox}
\linead{\linewidth}\\[1mm]
\linead{\linewidth}

\begin{titlebox}{milcVino}
Compromiso final
\end{titlebox}

\small
\begin{tabularx}{\textwidth}{>{\RaggedRight\arraybackslash}p{3.0cm}YYY}
\rowcolor{milcVino}\color{white}\bfseries Criterio & \color{white}\bfseries Lo logré & \color{white}\bfseries En proceso & \color{white}\bfseries Necesito apoyo\\
Comprensión del tema & Explico ideas centrales con mis palabras. & Reconozco ideas pero falta precisión. & Necesito repasar conceptos base.\\
Pensamiento computacional & Apliqué los 4 pilares al diseñar mi producto. & Apliqué algunos pilares. & Necesito releer la fase 2.\\
Aplicación y producto & Mi producto cumple los criterios y se puede revisar. & Producto funcional con ajustes pendientes. & Aún en construcción.\\
Reflexión 5 dimensiones & Respondí cada dimensión con honestidad. & Respondí algunas con detalle. & Mis respuestas son muy generales.\\
Triángulo de pensamiento & Conecté las 3 voces con mi trabajo real. & Conecté al menos una. & No relaciono las voces con mi proceso.\\
\end{tabularx}
\normalsize

\textbf{Hoy entendí que:}\\[1mm]
\linead{\linewidth}\\[1mm]
\linead{\linewidth}

\textbf{Mi compromiso para la próxima sesión es:}\\[1mm]
\linead{\linewidth}\\[1mm]
\linead{\linewidth}

\begin{softbox}{milcMostaza}{milcGris}{Cita de despedida · Marco Aurelio}
\textit{``El obstáculo en el camino se vuelve el camino. Lo que estorba al camino se vuelve camino.''}\\[1mm]
Cada error de hoy, bien observado, se convierte en pieza del camino que pisarás mañana.
\end{softbox}

\small
\begin{tabularx}{\textwidth}{>{\bfseries}p{3.6cm}Y}
\rowcolor{milcGradeDark!12}Nombre completo: & \linead{10cm}\\
Fecha: & \linead{4cm} \quad Curso/grupo: \linead{4cm}\\
Mi firma: & \linead{9cm}\\
\end{tabularx}
\normalsize

\begin{infoband}{milcGradeDark}
\textbf{Para la próxima sesión.} Llega con el pitch deck terminado, guion escrito y grabación corregida. En la sesión 7 vas a trabajar la \textbf{marca e identidad visual} del proyecto: nombre, paleta, logo, tono. La marca se construye sobre la narrativa del pitch; sin pitch claro hoy, la marca de la próxima sesión será adorno sin alma.
\end{infoband}

\end{document}
